
A correctly optimized score network $s_\chi(r,t)$ allows, through simulation of \Cref{eq:bw_sde_j}, to obtain samples from the joint distribution $q_\psi(z)$. Similarly, a \textit{conditional} score network $s_\chi(r^{A_1},t\g z^{A_2})$ allows, through the simulation of \Cref{eq:cond_bw_sde_c}, to sample from $q_\psi(z^{A_1}\g z^{A_2})$. In \Cref{sec:multitime} we extend guidance mechanisms used in classical diffusion models to allow multi-modal conditional generation. A na\"ive alternative is to rely on the unconditional score network $s_\chi(r,t)$ for the conditional generation task, by casting it as an \textit{in-painting} objective. Intuitively, any missing modality could be recovered in the same way as a uni-modal diffusion model can recover masked information.%, e.g. a black patch on an image. 
In \Cref{sec:inpainting} we discuss the implicit assumptions underlying in-painting from an information theoretic perspective, and argue that, in the context of multi-modal data, such assumptions are difficult to satisfy. Our intuition is corroborated by ample empirical evidence, where our method consistently outperform alternatives.

\subsection{Multi-time Diffusion}
\label{sec:multitime}

We propose a modification to the classifier-free guidance technique \citep{ho2022classifier} to learn a score network that can generate conditional and unconditional samples from any subset of modalities. Instead of training a separate score network for each possible combination of conditional modalities, which is computationally infeasible, we use a single architecture that accepts all modalities as inputs and a \textit{multi-time vector} $\tau=[t_1,\dots,t_M]$. The multi-time vector serves two purposes: it is both a conditioning signal and the time at which we observe the diffusion process. 

\noindent \textbf{Training:} learning the conditional score network relies on randomization. As discussed in \Cref{sec:mld}, we consider an arbitrary partitioning of all modalities in two disjoint sets, $A_1$ and $A_2$. The set $A_2$ contains randomly selected conditioning modalities, while the remaining modalities belong to set $A_1$. Then, during training, the parametric score network estimates $\nabla\log(q(r,t\g z^{A_2}))$, whereby the set $A_2$ is randomly chosen at every step. This is achieved by the \textit{masked diffusion process} from \Cref{eq:cond_fw_sde_c}, which only diffuses modalities in $A_1$. More formally, the score network input is $R_t=\mathcal{C}(R^{A_1}_t,Z^{A_2})$, along with a multi-time vector ${\tau}(A_1,t) = t\begin{bmatrix} \mathds{1}(1\in A_1),\dots,\mathds{1}(M\in A_1) \end{bmatrix}$.
As a follow-up of the example in \Cref{sec:mld}, given $A_1 = \{1,3,5\}$, such that $X^{A_1} = \{X^1, X^3, X^5 \}$, and $A_2 = \{2,4,6\}$ such that $X^{A_2} = \{X^2, X^4, X^6 \}$, then, ${\tau}(A_1,t) = [t,0,t,0,t,0]$.


More precisely, the algorithm for the multi-time diffusion training (see \ref{apdx:mld_details} for the pseudo-code) proceeds as follows. At each step, a set of conditioning modalities $A_2$ is sampled from a predefined distribution $\nu$, where $\nu(\emptyset)\defeq\Pr(A_2=\emptyset)=d$, and $\nu(U)\defeq\Pr(A_2=U)=\nicefrac{(1-d)}{(2^M-1)}$ with $U\in \mathcal{P}(\{1,\dots,M\})\setminus\emptyset$, where $\mathcal{P}(\{1,\dots,M\})$ is the powerset of all modalities. The corresponding set $A_1$ and mask $m(A_1)$ are constructed, and a sample $X$ is drawn from the training data-set. The corresponding latent variables $Z^{A_1}=\{e^i_{\psi}(X^i)\}_{i\in A_1}$ and $Z^{A_2}=\{e^i_{\psi}(X^i)\}_{i\in A_2}$ are computed using the pre-trained encoders, and a diffusion process starting from $R_0 = \mathcal{C}(Z^{A_1}, Z^{A_2})$ is simulated for a randomly chosen diffusion time $t$, using the conditional forward SDE with the mask $m(A_1)$. The score network is then fed the current state $R_t$ and multi-time vector $\tau(A_1,t)$, and the difference between the score network's prediction and the true score is computed, applying the mask $m(A_1)$. The score network parameters are updated using stochastic gradient descent, and this process is repeated for a total of $L$ training steps. Clearly, when $A_2 = \emptyset$, training proceeds as for an un-masked diffusion process, since the mask $m(A_1)$ allows all latent variables to be diffused.






\noindent \textbf{Conditional generation:} any valid numerical integration scheme for \Cref{eq:cond_bw_sde_c} can be used for conditional sampling (see \ref{apdx:mld_details} for an implementation using the Euler-Maruyama integrator). First, conditioning modalities in the set $A_2$ are encoded into the corresponding latent variables $z^{A_2}=\{e^j(x^{j})\}_{j\in A_2}$. Then, numerical integration is performed with step-size $\Delta t=\nicefrac{T}{N}$, starting from the initial conditions $R_0=\mathcal{C}(R^{A_1}_0,z^{A_2})$, with $R^{A_1}_0\sim \rho(r^{A_1})$. At each integration step, the score network $s_\chi$ is fed the current state of the process and the multi-time vector $\tau(A_1,\cdot)$. Before updating the state, the masking is applied. Finally, the generated modalities are obtained thanks to the decoders as $\hat{X}^{A_1}=\{d^j_\theta(R^{j}_T)\}_{j\in A_1}$.
Inference time conditional generation is not randomized: conditioning modalities are the ones that are available, whereas the remaining are the ones we wish to generate.

%Inference time conditional generation is not randomized: conditioning modalities are the ones that are available, whereas the remaining modalities are the ones we wish to generate.

%, given the available modalities. 


 Any-to-any multi-modality has been recently studied through the composition of modality-specific diffusion models \citep{tang2023anytoany}, by designing cross-attention and training procedures that allow arbitrary conditional generation. The work by \citet{tang2023anytoany} relies on latent interpolation of input modalities, which is akin to mixture models, and uses it as conditioning signal for individual diffusion models. This is substantially different from the joint nature of the multi-modal latent diffusion we present in our work: instead of forcing entanglement through cross-attention between score networks, our model relies on joint diffusion process, whereby modalities naturally co-evolve according to the diffusion process. Another recent work \citep{wu2023nextgpt} targets multi-modal conversational agents, whereby the strong, underlying assumption is to consider one modality, i.e., text, as a guide for the alignment and generation of other modalities. Even if conversational objectives are orthogonal to our work, techniques akin to instruction following for cross-generation, are an interesting illustration of the powerful capabilities of in-context learning of LLMs \citep{xie2022an, min2022rethinking}.


\subsection{In-painting and its implicit assumptions}\label{sec:inpainting}
Under certain assumptions, given an unconditional score network $s_\chi(r,t)$ that approximates the true score $\nabla\log q(r,t)$, it is possible to obtain a conditional score network $s_\chi(r^{A_1},t\g z^{A_2})$, to approximate $\nabla\log q(r^{A_1},t\g z^{A_2})$. We start by observing the equality:

\begin{footnotesize}
\begin{equation}\label{eq:inp}
    q(r^{A_1},t\g z^{A_2})=\int q(\mathcal{C}(r^{A_1},r^{A_2}),t\g z^{A_2}) \d r^{A_2}=\int \frac{q(z^{A_2}\g \mathcal{C}(r^{A_1},r^{A_2}),t)}{q_\psi(z^{A_2})}q(\mathcal{C}(r^{A_1},r^{A_2}),t)\d r^{A_2},
\end{equation}
\end{footnotesize}


where, with a slight abuse of notation, we indicate with $q(z^{A_2}\g \mathcal{C}(r^{A_1},r^{A_2}),t)$ the density associated to the event: the portion corresponding to $A_2$ of the latent variable $Z$ is equal to $z^{A_2}$ given that the whole diffused latent $R_t$ at time $t$, is equal to $\mathcal{C}(r^{A_1},r^{A_2})$. In the literature, the quantity $q(z^{A_2}\g \mathcal{C}(r^{A_1},r^{A_2}),t)$ is typically approximated by dropping its dependency on $r^{A_1}$. This approximation can be used to manipulate \Cref{eq:inp} as $q(r^{A_1},t\g z^{A_2})\simeq \int q(r^{A_2},t\g z^{A_2})q(r^{A_1},t|r^{A_2},t)\d r$. Further Monte-Carlo approximations \citep{Song2020,lugmayr2022repaint} of the integral allow implementation of a practical scheme, where an approximate conditional score network is used to generate conditional samples. This approach, known in the literature as \textit{in-painting}, provides high quality results in several \textit{uni-modal} application domains \citep{Song2020,lugmayr2022repaint}. 
% \noteGF{Let's discuss live!}Intuitively, the conditional probability approximation can be justified as follows: to reconstruct a patch of an image (in our notation, that would be $r^{A_1}$), given observation of a noisy version of the patch or, observing a noisy version of the rest of the image (in our notation, that would be $z^{A_2}$), makes limited to no difference. 

The \gls{KL} divergence between $q(z^{A_2}\g \mathcal{C}(r^{A_1},r^{A_2}),t)$ and $q(z^{A_2}\g r^{A_2},t)$ quantifies, fixing $r^{A_1},r^{A_2}$, the discrepancy between the true and approximated conditional probabilities. Similarly, the expected \gls{KL} divergence $\Delta=\int q(r,t)\mathrm{KL}[q(z^{A_2}\g \mathcal{C}(r^{A_1},r^{A_2}),t)\g\g q(z^{A_2}\g r^{A_2},t)]\dd r$, provides information about the average discrepancy. Simple manipulations allow to recast this as a discrepancy in terms of mutual information $\Delta=I(Z^{A_2};R^{A_1}_t,R^{A_2}_t)-I(Z^{A_2};R^{A_2}_t)$. Information about $Z^{A_2}$ is contained in $R^{A_2}_t$, as the latter is the result of a diffusion with the former as initial conditions, corresponding to the Markov chain $R^{A_2}_t\rightarrow Z^{A_2}$, and in $R^{A_1}_t$ through the Markov chain $Z^{A_2} \rightarrow Z^{A_1} \rightarrow R^{A_1}_t$. The positive quantity $\Delta$ is close to zero whenever the rate of loss of information w.r.t initial conditions is similar for the two subsets $A_1,A_2$. In other terms, $\Delta\simeq 0$ whenever out of the whole $R_t$, the portion $R^{A_2}_t$ is a sufficient statistic for $Z^{A_2}$.  

The assumptions underlying the approximation are in general not valid in the case of multi-modal learning, where the robustness to stochastic perturbations of latent variables corresponding to the various modalities can vary greatly. Our claim are supported empirically by an ample analysis on real data in \ref{apdx:mld_variants}, where we show that multi-time diffusion approach consistently outperforms in-painting.

